\scanpage{193}
\mspage{499}

le sur la vie de mes parents et sur mon enfance. C'est des mois ou des années pl\unsure{us} tard seulement, je crois, que par un certain nombre de recoupements, je me suis rendu compte que dans certains cas mon appréhension des rôles yin et yang dans tels ou tels ``couples'' était resté un tantinet superficielle; que j'avais mis da\unsure{ns} le même sac, un peu hâtivement, des situations de nature différente que la diale\unsure{c}tique yin-yang chinoise prenait bien soin de distinguer (112\handwritten{$'$}). Maintenant, je m\unsure{e} rends compte que l'appréhension du yin et du yang reste encore relativement grossière et statique chez moi, surtout si on la compare à la finesse requise pour l'exercice de certains arts traditionnels chinois comme la médecine (intimement liée aussi à la diététique et à l'art culinaire), où cette appréhension finit pa\unsure{r} devenir comme une seconde nature.

J'ai eu l'impression plus d'une fois que chez les pratiquants et praticiens de ces arts, qu'ils soient orientaux ou européens, cette finesse d'appréhension reste fragmentaire, en ce sens qu'elle reste, dans une très large mesure, soigneusement cantonnée dans l'exercice de cet art. Dans la vie de tous les jours, elle agirait plutôt comme un ``savoir'' ordinaire, se superposant purement et simplement au ``savoir'' du conditionnement culturel (et autre), et restant plus ou moins lettre morte vis-à-vis de celui-ci. Pour le dire autrement, j'ai eu l'impression que la vision du monde et de soi, et les mécanismes de répression dans la perception de la réalité, ne sont en rien différents chez ces personnes tout ce qu'il y a d'``averties'', que chez le commun des mortels.

Cette impression se recoupe avec une autre, que j'ai eue en parcourant deux ou trois textes, écrits par des européens censés être ``dans le coup'', qui se proposent de donner un aperçu de la philosophie traditionelle chinoise du yin et du yang. (L'un des auteurs est un orientaliste français bien connu, dont le nom m'échappe à présent.) La chose qui m'a frappée, c'est que dans ces textes, le yin et le yang sont présentés comme des principes ``o p p o s é s'' (ou ``c o n t r a i r e s'') voire a n t a g o n i s t e s (ce dernier terme revient à plusieurs reprises dans un de ces textes), plutôt que c o m p l é m e n t a i r e s. Cette ``opposition'' ou ``antagonime'' aurait son expression typique dans celui qui aurait lieu entre la femme et l'homme à l'intérieur de la société humaine, et à l'intérieur du couple institué par la société.

L'antagonisme dans le couple époux-épouse est bel et bien une réalité, aussi bien à l'Est qu'à l'Ouest. Elle est profondément enracinée dans la culture, à tel point qu'il peut paraître parfois comme un des aspects (parfois déroutant!) de la condition humaine, voire même comme la racine du conflit en l'homme ou dans la société humaine. La réalité de cet antagonisme est irrécusable, et elle dépasse certes les clichés courants qui s'efforcent de l'exorciser

% archive sheet number: 192
